\chapter{Singular Reduced Hessians}
\label{chap:singular}

\vspace*{\fill}

The nullspace method requires the reduced Hessian to be positive definite
(\cref{ass:nullspace}~(ii)). The Hessian of the regularized master problem is only positive
semidefinite, and its rank is bounded by the number of linking rows, while the number of
columns keeps growing. It is therefore unavoidable that a reduction of the working set
occasionally produces a nullspace direction along which the objective has no curvature.

\medskip

This chapter shows where such singularities can arise, how they are detected from the
pivots of the factorization, and how they are resolved. The resolution is a
\textbf{zero-curvature step}: along a direction in the kernel of the reduced Hessian, the
objective decreases linearly, and the step is taken until a variable reaches its bound. We
prove that adding this variable to the working set removes the singularity, so that the
interrupted reduction can be completed. The treatment covers released columns and released
surpluses with one argument, and keeps every working-set variable exactly at its bound.

\vspace*{\fill}

\newpage


\section{Where Singularity Can Arise}
\label{sec:where-singular}

By \cref{cor:dimension}, the reduced Hessian is positive definite if and only if
$\Heff = (\Aeff Z_x)^\top(\Aeff Z_x)$ is, i.e.\ if and only if the vectors
$\Aeff z_j$, $j\in\Nx$, are linearly independent. \Cref{prop:expand-pd} shows that expansions
preserve this property. Singularity can therefore only arise in one of the two reductions
of \cref{thm:H-reduce}.

\begin{prop}[Singular reductions] \label{prop:singular-cases}
Let $\Heff$ be positive definite before the reduction.
\begin{enumerate}
    \item[(iv)] After releasing column $r$, the matrix $\Heff'$ of~\eqref{eq:bordered} is
          singular if and only if $\Aeff z_r\in\spn\{\Aeff z_k : k\in\Nx\}$, which is the case if
          and only if the Schur complement $\sigma_r = h_{rr} - h^\top\Heff^{-1}h$ vanishes.
    \item[(v)] After releasing the surplus of row $i$, the matrix $\Heff' = \Heff - v_iv_i^\top$
          is singular if and only if $\omega_i = 1 - v_i^\top\Heff^{-1}v_i$ vanishes.
\end{enumerate}
Moreover $\sigma_r\ge0$ and $0\le\omega_i\le1$.
\end{prop}

\begin{proof}
(iv) $\Heff'$ is the Gram matrix of the vectors $\Aeff z_k$, $k\in\Nx$, and $u = \Aeff z_r$. It
is singular if and only if these vectors are linearly dependent; since the vectors $\Aeff z_k$
alone are independent, this means $u\in\spn\{\Aeff z_k\}$. For the bordered matrix
$\Heff' = \left[\begin{smallmatrix}\Heff & h\\ h^\top & h_{rr}\end{smallmatrix}\right]$, block
elimination gives $\det\Heff' = \det\Heff\cdot(h_{rr} - h^\top\Heff^{-1}h) = \det\Heff\cdot\sigma_r$,
so $\Heff'$ is singular if and only if $\sigma_r = 0$.

(v) By the matrix determinant lemma, $\det(\Heff - v_iv_i^\top) = \det\Heff\cdot(1 - v_i^\top\Heff^{-1}v_i)
= \det\Heff\cdot\omega_i$, so $\Heff'$ is singular if and only if $\omega_i = 0$.

In both cases $\Heff'$ is a Gram matrix, of $[\,\Aeff Z_x\;\;u\,]$ in case~(iv) and of $\Aeff'Z_x$
($\Aeff$ without row $i$) in case~(v). It is therefore positive semidefinite, so
$\det\Heff'\ge0$, and with $\det\Heff>0$ we get $\sigma_r\ge0$ and $\omega_i\ge0$. Finally
$\omega_i\le1$ because $\Heff^{-1}$ is positive definite, so $v_i^\top\Heff^{-1}v_i\ge0$.
\end{proof}

\Cref{prop:singular-cases} also shows that a downdate in case~(v) can never make $\Heff$
indefinite: in exact arithmetic, the worst case is a singular matrix.


\section{Detection}
\label{sec:detection}

Testing $\sigma_r$ or $\omega_i$ before the update would cost an additional solve for every
reduction. Instead, the update is performed, and the pivots $D_{kk}$ of the factorization
$\Heff^{\mathrm{emb}} = LDL^\top$ are inspected afterwards.

\begin{definition}[Pivot invariant] \label{def:pivot-invariant}
For $k\in\Nx$, let $\gamma_k = \lVert a_k - a_{\rep(k)}\rVert^2$ be the \textbf{pivot scale} of
index $k$ (with all rows, not only those of $\Aeff$). The factorization satisfies the
\textbf{pivot invariant} with tolerance $\varepsilon_{\mathrm{piv}}>0$ if
\[
D_{kk}\;>\;\varepsilon_{\mathrm{piv}}\,\gamma_k\qquad\text{for all }k\in\Nx .
\]
\end{definition}

The scale $\gamma_k$ uses all rows on purpose: a column whose coefficients lie entirely in rows
with free surpluses has $(\Heff)_{kk} = 0$, and a scale based on $\Heff$ would hide this
singularity.

The determinant of $\Heff'$ is the product of its pivots. In case~(iv) the update adds one new
pivot, and in case~(v) the pivots change by the factor $\omega_i$ in total, since
$\det\Heff' = \det\Heff\cdot\omega_i$. If all pivots satisfied the invariant before the update, a
vanishing $\sigma_r$ or $\omega_i$ must therefore show up as a pivot that violates it after the
update. This detection is a numerical heuristic rather than a theorem about floating-point
arithmetic, but it follows directly from the determinant identities of
\cref{prop:singular-cases}. When it fires, the update is undone by a refactorization with the
previous working set (\cref{def:refactorization}), which also removes any large entries
created by dividing by a tiny pivot. We use $\varepsilon_{\mathrm{piv}} = 10^{-10}$.


\section{Zero-Curvature Directions}
\label{sec:zero-curvature}

Assume now that the iterate $y$ is a subspace minimizer for the working set $\Wset$, that the
index to be released has a negative multiplier, and that its release would make the reduced
Hessian singular. We construct a direction along which the objective decreases linearly and
the working set (apart from the released index) is respected. Recall from
\cref{cor:multipliers-minimizer} what is known at a subspace minimizer:
\begin{equation} \label{eq:minimizer-facts}
\begin{aligned}
g^\top z_k &= g_k - g_{\rep(k)} = 0 \;\;(k\in\Nx), &\qquad \pi_\ell &= 0 \;\;(\ell\in\Ns),\\
g_r - g_{\rep(r)} &= (\nu_x)_r \;\;(r\in\Wx), & \pi_i &= (\nu_s)_i \;\;(i\in\Ws).
\end{aligned}
\end{equation}
In this section $z_k = e_k - e_{\rep(k)}$ denotes the column part in $\R^n$.

\subsection{Released column}

\begin{theorem}[Zero-curvature direction, column] \label{thm:zc-column}
Let $r\in\Wx$ have the multiplier $(\nu_x)_r<0$, and suppose that the release of $r$ is
singular, i.e.\ $\Aeff z_r = \sum_{k\in\Nx}t_k\,\Aeff z_k$ for some coefficients $t$ (these are
$t = \Heff^{-1}h$ with $h$ from~\eqref{eq:bordered}). Define $d = [d_x; d_s]$ by
\begin{equation} \label{eq:zc-column}
d_x = z_r - \sum_{k\in\Nx}t_k\,z_k,\qquad
(d_s)_\ell = (A_Id_x)_\ell\;\;(\ell\in\Ns),\qquad (d_s)_\ell = 0\;\;(\ell\in\Ws).
\end{equation}
Then
\begin{enumerate}
    \item $Cd_x = 0$, and $d_\ell = 0$ for every $\ell\in\Wset\setminus\{r\}$;
    \item $\Abar d = 0$;
    \item $\nabla f(y)^\top d = (\nu_x)_r < 0$.
\end{enumerate}
\end{theorem}

\begin{proof}
(i) Every $z_k = e_k - e_{\rep(k)}$ has one entry $+1$ and one entry $-1$ in the same block, so
$Cz_k = 0$ and therefore $Cd_x = 0$. The vector $d_x$ is supported on $\{r\}\cup\Nx\cup\Sset$, so it
vanishes on $\Wx\setminus\{r\}$, and $d_s$ vanishes on $\Ws$ by definition.

(ii) The rows of $\Abar d$ are $A_Ed_x$ (equality rows) and $A_Id_x - d_s$ (inequality rows). On
the rows of $\Aeff$, i.e.\ the equality rows and the inequality rows in $\Ws$ (where $d_s = 0$),
$\Abar d$ equals $\Aeff d_x = \Aeff z_r - \sum_kt_k\Aeff z_k$, which is zero by the assumption on
$t$. On an inequality row $\ell\in\Ns$, $(\Abar d)_\ell = (A_Id_x)_\ell - (d_s)_\ell = 0$ by the
definition of $d_s$.

(iii) Splitting $g^\top d = g_x^\top d_x + \sum_\ell\pi_\ell(d_s)_\ell$ (recall $g_s = \pi_I$) and
using~\eqref{eq:minimizer-facts},
\[
g^\top d = \bigl(g_r - g_{\rep(r)}\bigr) - \sum_{k\in\Nx}t_k\underbrace{\bigl(g_k - g_{\rep(k)}\bigr)}_{=0}
+ \sum_{\ell\in\Ns}\underbrace{\pi_\ell}_{=0}(d_s)_\ell = (\nu_x)_r . \qedhere
\]
\end{proof}

\subsection{Released surplus}

\begin{theorem}[Zero-curvature direction, surplus] \label{thm:zc-surplus}
Let the surplus of row $i\in\Ws$ have the multiplier $(\nu_s)_i = \pi_i < 0$, and suppose that
its release is singular, $\omega_i = 1 - v_i^\top\Heff^{-1}v_i = 0$. Let $t = \Heff^{-1}v_i$, so
that $v_i^\top t = 1$. Define $d = [d_x;d_s]$ by
\begin{equation} \label{eq:zc-surplus}
d_x = \sum_{k\in\Nx}t_k\,z_k,\qquad
(d_s)_\ell = (A_Id_x)_\ell\;\;(\ell\in\Ns\cup\{i\}),\qquad (d_s)_\ell = 0\;\;\text{otherwise}.
\end{equation}
Then
\begin{enumerate}
    \item $Cd_x = 0$, $d_x$ vanishes on $\Wx$, $d_s$ vanishes on $\Ws\setminus\{i\}$, and $(d_s)_i = 1$;
    \item $\Abar d = 0$;
    \item $\nabla f(y)^\top d = \pi_i < 0$.
\end{enumerate}
\end{theorem}

\begin{proof}
(i) As in \cref{thm:zc-column}, $Cd_x = 0$ and $d_x$ vanishes on $\Wx$. Further, $d_x = Z_xt$,
so $(d_s)_i = (A_IZ_xt)_i = (A_I)_{i,:}Z_x\,t = v_i^\top t = 1$, using
$v_i = Z_x^\top(A_I)_{i,:}^\top$ from~\eqref{eq:rank-one-up}.

(ii) Let $\Aeff'$ be $\Aeff$ without row $i$; its rows are the equality rows and the inequality
rows in $\Ws\setminus\{i\}$. By \cref{thm:H-reduce}, $(\Aeff'Z_x)^\top(\Aeff'Z_x) = \Heff - v_iv_i^\top$,
hence
\[
\lVert\Aeff'd_x\rVert^2 = t^\top(\Heff - v_iv_i^\top)\,t = t^\top v_i - (v_i^\top t)^2 = 1 - 1 = 0,
\]
using $\Heff t = v_i$ and $v_i^\top t = 1$. So $\Aeff'd_x = 0$, which means that $\Abar d$ vanishes
on the rows of $\Aeff'$ (where $d_s = 0$). On the remaining inequality rows
$\ell\in\Ns\cup\{i\}$, $(\Abar d)_\ell = (A_Id_x)_\ell - (d_s)_\ell = 0$ by construction.

(iii) Using~\eqref{eq:minimizer-facts}, $g_x^\top d_x = \sum_kt_k(g_k - g_{\rep(k)}) = 0$ and
$\pi_\ell = 0$ for $\ell\in\Ns$, so only the component of surplus $i$ remains:
$g^\top d = \pi_i\,(d_s)_i = \pi_i$.
\end{proof}

\begin{remark} \label{rem:zc-kappa}
For a ``$\le$'' row with slack coefficient $\kappa_i = +1$ (\cref{rem:kappa-summary}), the
column part is scaled by $-\kappa_i$, $d_x = -\kappa_i\sum_kt_kz_k$, and the surplus part is
$(d_s)_\ell = -\kappa_\ell(A_Id_x)_\ell$. Then again $(d_s)_i = 1$ and
$g^\top d = -\kappa_i\pi_i = (\nu_s)_i < 0$.
\end{remark}

\begin{cor}[Linear decrease] \label{cor:linear-decrease}
For the directions of \cref{thm:zc-column,thm:zc-surplus},
\[
f(y + \alpha d) = f(y) + \alpha\,\nabla f(y)^\top d,\qquad \pi(y+\alpha d) = \pi(y)
\qquad\text{for all }\alpha .
\]
\end{cor}

\begin{proof}
Expanding the quadratic $f$, $f(y+\alpha d) = f(y) + \alpha\nabla f(y)^\top d +
\frac{\alpha^2}{2\delta}\lVert\Abar d\rVert^2$, and $\Abar d = 0$. The dual estimate
$\pi(y) = \pihat + \frac{1}{\delta}(b - \Abar y)$ depends on $y$ only through $\Abar y$, which does
not change.
\end{proof}

\begin{remark}[Fully masked columns] \label{rem:masked-column}
If all nonzero coefficients of $u = a_r - a_{\rep(r)}$ lie in rows with free surpluses, then
$\Aeff z_r = 0$, so $h = 0$ and $h_{rr} = 0$. \Cref{thm:zc-column} then applies with $t = 0$:
the direction is $z_r$ together with the surplus components that absorb $A_Iz_r$.
\end{remark}


\section{The Zero-Curvature Step Removes the Singularity}
\label{sec:zc-step}

Along $d$ the objective decreases linearly (\cref{cor:linear-decrease}), so the method moves
as far as the bounds allow. In this section $\rho$ denotes the released index: the column $r$ of
\cref{thm:zc-column} or the surplus $i$ of \cref{thm:zc-surplus}.

\begin{lemma}[Finite step] \label{lem:finite-step}
The directions of \cref{thm:zc-column,thm:zc-surplus} have at least one negative component.
Hence the ratio test
\[
\alpha^\star = \min_{\ell:\;d_\ell<0}\;-\frac{y_\ell}{d_\ell}
\]
is finite, and it is attained at a \textbf{blocking index} $q$ in
$\Sset\cup\Nx\cup\Ns$.
\end{lemma}

\begin{proof}
Suppose $d\ge0$. By part~(i) of the theorems, every block sum of $d_x$ is zero ($Cd_x = 0$).
A nonnegative vector whose block sums are zero is zero, so $d_x = 0$. Then
$\Abar d = [\,A_Ed_x;\; A_Id_x - d_s\,] = [\,0;\;-d_s\,]$, and $\Abar d = 0$ (part~(ii)) forces
$d_s = 0$. But $d = 0$ contradicts $\nabla f(y)^\top d < 0$ (part~(iii)). So some component of
$d$ is negative, and the minimum is taken over a nonempty set. The blocking index cannot lie in
$\Wset\setminus\{\rho\}$, where $d$ vanishes, and it is not the released index $\rho$, whose
component of $d$ is positive ($d_r = 1$ or $(d_s)_i = 1$). Hence $q\in\Sset\cup\Nx\cup\Ns$.
\end{proof}

After the step, the blocking index $q$ is added to the working set by an expansion
(\cref{sec:expanding}), and the release is attempted again. The following theorem shows that
this attempt cannot fail for the same reason.

\begin{theorem}[Removal of the dependency] \label{thm:removal}
Let $\rho$ be the released index and $q$ the blocking index of
\cref{lem:finite-step}. Then the reduced Hessian of the working set
$\Wset' = (\Wset\setminus\{\rho\})\cup\{q\}$ is positive definite. In other words, after the
zero-curvature step and the expansion by $q$, the retried release of $\rho$ succeeds.
\end{theorem}

\begin{proof}
For any working set with nullspace basis $Z$, the reduced Hessian $Z^\top\Abar^\top\Abar Z$ is
positive definite if and only if $\Abar Z\eta\ne0$ for all $\eta\ne0$, i.e.\ if and only if
$\Abar$ is injective on the nullspace $\Null(D)$ of the working set. We use this for three
spaces:
\begin{align*}
N_0 &= \Null(D_\Wset) && \text{(before the release)},\\
V &= \Null(D_{\Wset\setminus\{\rho\}}) && \text{(after the release)},\\
V' &= V\cap\{y : y_q = 0\} = \Null(D_{\Wset'}) && \text{(after the release and the expansion by $q$)}.
\end{align*}
\emph{Step 1: $\ker\Abar\cap V = \spn\{d\}$.} Before the release the reduced Hessian was positive
definite, so $\ker\Abar\cap N_0 = \{0\}$. Removing one constraint from the working set adds one
dimension, $N_0\subset V$ with $\dim V = \dim N_0 + 1$. If $\ker\Abar\cap V$ contained two
linearly independent vectors, some nonzero combination of them would satisfy the one extra
constraint ($y_\rho = 0$) and therefore lie in $N_0\cap\ker\Abar = \{0\}$, a contradiction. So
$\dim(\ker\Abar\cap V)\le1$. The direction $d$ lies in $V$ (part~(i) of
\cref{thm:zc-column,thm:zc-surplus}) and in $\ker\Abar$ (part~(ii)), and $d\ne0$, so
$\ker\Abar\cap V = \spn\{d\}$.

\emph{Step 2: $\ker\Abar\cap V' = \{0\}$.} Every vector of $\ker\Abar\cap V'$ is a multiple
$\beta d$ of $d$ by Step~1, and it has a zero component $q$. Since $d_q<0$, this forces
$\beta = 0$. Hence $\Abar$ is injective on $V' = \Null(D_{\Wset'})$, and the reduced Hessian of
$\Wset'$ is positive definite.
\end{proof}

The retried release starts from the same kind of state as the first attempt: a subspace
minimizer with a negative multiplier for $\rho$.

\begin{prop}[The iterate stays a subspace minimizer] \label{prop:stays-minimizer}
After the zero-curvature step and the expansion by $q$, the iterate is a subspace minimizer
of the new working set $(\Wset\setminus\{\rho\})\cup\{q\}$ with $\rho$ still fixed, and the
multiplier of $\rho$ is unchanged.
\end{prop}

\begin{proof}
By \cref{cor:linear-decrease}, $\pi$ and hence the gradient $g = \cbar - \Abar^\top\pi$ do not change
along the step. The expansion by $q$ changes the nullspace basis as in \cref{thm:Z-update}: it
removes a vector, or, if $q$ is a representative, replaces vectors $z_j$ by
$z_j - z_{r'}$ for the replacement $r'$. In all cases the new basis vectors are combinations of
old ones, whose reduced gradient entries $g^\top z$ were zero; so the new reduced gradient is zero
as well. The multiplier of $\rho$ is computed from the unchanged gradient by
\cref{cor:multipliers-minimizer}.
\end{proof}

\begin{remark}[The representative of the released column blocks] \label{rem:rep-blocks}
In \cref{thm:zc-column}, the blocking index may be the representative $\rep(r)$ of the block
of $r$. Then $r$, which now has the positive weight $\alpha^\star$, can simply take over as
the new representative (case~(ii) of \cref{lem:expand-ordering} with replacement $r$). This
completes the release at once.
\end{remark}

\Cref{alg:singular} summarizes the procedure. Each pass either completes the release or adds
one more index to the working set. The number of passes is therefore bounded by the number of
free variables.

\begin{algobox}{Release with Singularity Resolution}{alg:singular}
\State \textbf{Input:} subspace minimizer $y$, index $\rho$ (column $r$ or surplus $i$) with
       negative multiplier
\Loop
    \State Apply the reduction update (\code{rowadd} for a column, downdate for a surplus)
    \If{the pivot invariant holds}
        \State \textbf{return} \Comment{release completed}
    \EndIf
    \State Undo by refactorization with $\rho$ still in the working set
    \State Solve $\Heff t = h$ (column) or $\Heff t = v_i$ (surplus)
    \State Form $d$ by~\eqref{eq:zc-column} or~\eqref{eq:zc-surplus}
    \State Ratio test along $d$; $y\gets y+\alpha^\star d$ \Comment{$\pi$ is unchanged}
    \If{$\rho = r$ and $q = \rep(r)$}
        \State Make $r$ the representative, put $q$ into $\Wx$; \textbf{return}
               \Comment{\cref{rem:rep-blocks}}
    \EndIf
    \State Expand the working set by $q$ \Comment{\cref{thm:removal}: the next pass succeeds}
\EndLoop
\end{algobox}

\begin{remark}[Why not remove the largest coefficient?] \label{rem:argmax}
A simpler rule moves the index $k = \arg\max_{k\in\Nx}|t_k|$ of the largest coefficient of the
linear dependence into the working set and retries. This also removes the dependency, but
it ignores the value of $x_k$. If $x_k>0$, the working-set constraint $x_k = 0$ is violated by
the iterate, and its multiplier no longer certifies optimality: the method can stop with
$x_k>0$ and a positive multiplier. The zero-curvature step avoids this. The iterate stays
feasible, every working-set variable is exactly at its bound, the objective decreases
strictly, and the index added to the working set is the one selected by a ratio test rather
than by a heuristic. The same argument covers released columns and released surpluses.
\end{remark}
